% !TeX root = ../main.tex
\section{String and fret assignment (fingering)}
\label{sec:fingering}

\subsection{What the stage must do}

Input: notes (onset, pitch, duration), grouped into chords, plus a tuning.
Output: one string and fret per note, such that every chord is playable (one
note per string, a reachable fret span) and the sequence is comfortable and
idiomatic. Section~\ref{sec:problem} showed why this is ambiguous: a pitch has
up to five positions, a three-note chord can have dozens of shapes.

\subsection{Search with hand-designed costs}

\begin{keyidea}[: fingering as a shortest path]
Treat every possible position (or chord shape) of every note as a node, give
each node a cost (high frets, stretched shapes) and each move between
consecutive nodes a cost (hand movement), and find the cheapest path through
the whole piece with the Viterbi algorithm (Section~\ref{sec:viterbi}). Sayegh
framed it this way in 1989~\cite{sayegh1989fingering}; HMM versions decode
left-hand forms~\cite{hori2016minimax}.
\end{keyidea}

The weakness is that the costs are local and hand-made. They know nothing of
long-range habits: a riff repeated in another position, idiomatic chord shapes,
an open string used as a pedal tone, or the way solos move up the neck.

\subsection{Learned fingering}

\paragraph{Symbolic models} learn from human tabs (no audio). They treat
fingering as translation from note tokens to string/fret tokens, usually
pre-trained on DadaGP. MIDI-to-Tab~\cite{edwards2024miditotab} (encoder-decoder
transformer, masked pre-training, fine-tuned on professional transcriptions)
was rated more playable than other algorithms in a user study \reported.
Fretting-Transformer~\cite{hamberger2025fretting} (T5, with tuning and capo
conditioning) beat an A* baseline and Guitar Pro's own conversion \reported.
Hsu et al.~\cite{hsu2026noteevent} added note-event tokens and \term{constrained
decoding} that forbids string-fret pairs inconsistent with the pitch, and beat
Fretting-Transformer, especially with little training data \reported. None of
these three has released weights.

\paragraph{Audio-conditioned models} also listen: the same pitch sounds
different on a thick and a thin string. TART's third stage is a T5
conditioned on audio features~\cite{gupta2026tart}; with its note stage it
reaches 71.8\% string-fret Tab F1 averaged over its benchmarks, 8.5 points above
the best earlier baseline \reported. No code was found.

\begin{inapp}[: the Viterbi (fallback) \done]
\file{transcribe/arrange.rs}, \code{fret_events}. States are fingerings of each
chord (or single note); the hand position is carried through open strings. The
cost of a fingering is a weighted sum (\code{FingeringCosts}, weights fitted
with \code{eval_fingering} on GuitarSet players 00--02 and 10 human Guitar Pro
tabs):
\begin{center}
\small
\begin{tabular}{@{}lrl@{}}
\toprule
Term & Weight & Per \\
\midrule
\code{span} & 0.40 & fret between lowest and highest fretted note of a chord \\
\code{position} & $-0.01$ & fret of hand position (almost neutral after fitting) \\
\code{high} & 0.08 & fret above the 12th \\
\code{open} & 0.06 & open string \\
\code{gap} & 0.80 & unused string inside a chord (must be muted) \\
\code{moving} & 0.28 & fret of hand movement between events \\
\code{prior} & 0.11 & $-\ln P(\text{string}\mid\text{pitch})$, learned from 34\,150 GuitarSet notes \\
\bottomrule
\end{tabular}
\end{center}
Two hard-ish rules come from the techniques stage: a slide or hammer-on that
changes string costs 4.0 (\code{LINK_STRING_CHANGE}), a bend on an open string
6.0 (\code{BENT_OPEN_STRING}). Before fingering, a tuning is chosen from the
lowest notes (standard, E$\flat$, drop D, drop C, 7 and 8 strings).
\end{inapp}

\begin{inapp}[: the learned tagger \done{} (local model)]
\label{sec:fingering-app}
\file{transcribe/fingering_net.rs}, trained by \file{scripts/train/fingering/}.
\begin{itemize}
\item \textbf{Model}: a 4-layer Transformer \emph{tagger} (0.57\,M parameters).
  It does not generate tokens; for each note it outputs a score for each of
  the 6 strings. Inputs per note: the fret the note would have on each string
  (or ``impossible''), plus 16 timing and interval features (\code{common.features}).
  Windows of 128 notes.
\item \textbf{Decoder}: puts each chord on distinct strings, adds a stretch
  cost, and runs a hand-position Viterbi over the tagger's scores, with the
  same technique rules as above.
\item \textbf{Training data} (``ship'' variant, about 70\,k notes): GuitarSet
  00--02, EGDB train, Guitar-TECHS P1/P2. Settings were chosen with GuitarSet
  player 02 held out; never on 03--05.
\item \textbf{Use}: 6-string guitars, when \code{fingering_tagger.onnx} is in the
  models folder and the Transcribe form (or \cmd{--fingering learned}) asks for
  it; otherwise, and on any model error, the Viterbi. Local only (trained on
  EGDB).
\item \textbf{Parity}: the Rust port agrees with the Python model on 99.94\% of
  41\,293 test notes \measured. It needs each note's own onset: giving every
  chord note the chord's onset costs 0.02 accuracy.
\item \textbf{Cost}: about 6--10\ms{} per 128-note window on the CPU.
\end{itemize}
\end{inapp}

\subsection{What we measured}

\begin{table}[h]
\centering
\caption{String accuracy given the true notes, all test items, mean of 3 seeds
for the tagger \measured{} (\code{scripts/train/fingering/}). ``Full'' adds GOAT
training data (non-commercial).}
\label{tab:fingering}
\small
\begin{tabular}{@{}lccccc@{}}
\toprule
Method & GuitarSet comp & GuitarSet solo & EGDB test & EGSet12 & GOAT test \\
\midrule
Lowest fret (naive) & 0.33 & 0.25 & 0.65 & 0.43 & -- \\
Viterbi (app fallback) & 0.807 & 0.549 & 0.677 & 0.688 & 0.648 \\
Tagger, ship (app) & \textbf{0.853} & \textbf{0.616} & \textbf{0.806} & 0.698 & 0.667 \\
Tagger, full & 0.845 & 0.613 & 0.783 & \textbf{0.722} & \textbf{0.674} \\
\bottomrule
\end{tabular}
\end{table}

With transcribed notes (v2 CRNN) the tagger raises Tab F1 from 0.53--0.60 to
0.64--0.68 on the EGDB tones, from 0.53 to 0.63 on GuitarSet solo, and from
0.59 to 0.65 on BS-RoFormer stems of band mixes \measured. In the app pipeline
(\code{eval_guitarset}, players 03--05), the median share of notes on the
reference string and fret goes from 68\% to 78\% with GAPS notes on clean
guitar, and from 67\% to 76\% with v2 on ``Metal High Gain'' \measured.
Chords stay playable: fewer than 0.5\% impossible stretches, except 6.5\% on
EGSet12.

\subsection{Failure modes}
\begin{itemize}
\item \textbf{Too few players.} The tagger scores 0.95 and more on its training
  players and 0.57--0.69 on new ones; a 3.3\,M-parameter version did no better.
  Player diversity, not model size, is the limit \measured. DadaGP
  pre-training is the planned fix.
\item \textbf{Transposition hurts.} Training on passages shifted by $k$ frets
  lowered solo accuracy to 0.43: positions depend on the key \measured.
\item \textbf{Solos.} The Viterbi puts only 55--56\% of solo notes on the
  player's string (80\% for accompaniment): it keeps notes low while solos move
  up the neck.
\item \textbf{Licence.} The tagger is trained on EGDB, so it cannot be
  distributed until EGDB's authors allow it or it is retrained without EGDB.
\end{itemize}
